\documentclass[review]{elsarticle}

\usepackage[utf8]{inputenc}
\usepackage{amsmath,amssymb}
\usepackage{graphicx,booktabs,multirow,algorithm,algorithmic}
\usepackage{hyperref,xcolor}
\usepackage{float}

\journal{Renewable Energy}

\begin{document}

\begin{frontmatter}

\title{Liquid Neural Network-Regulated State Space Models for Probabilistic Wind Power Forecasting with NWP Integration}

\author{LNN-Mamba-Transformer}

\begin{abstract}
Probabilistic wind power forecasting is essential for reliable power system operation, yet existing approaches struggle to balance long-range temporal modeling with adaptive responses to non-stationary wind regimes. We propose LNMamba, a novel architecture that integrates Liquid Neural Network (LNN) dynamic gating with Liquid-Gated Selective State Space state space models (SSM) for 99-quantile probabilistic wind forecasting. The LNN gates produce input-dependent modulation signals that dynamically regulate the SSM's state transitions, enabling the model to adapt its temporal response to changing wind conditions without increasing parameter count. We evaluate LNMamba on two comprehensive NWP-enhanced datasets---GEFCom2012 (7 wind farms, ECMWF forecasts) and GEFCom2014 (10 zones)---totaling over 60,000 training windows across 17 sites. On GEFCom2012 Farm~1, LNMamba achieves a pinball loss of 0.0806, representing a 46.3\% improvement over the persistence baseline. Diebold-Mariano tests with Newey-West HAC standard errors confirm statistical significance at $p < 0.0001$ (DM = 12.432, 23/24 horizons individually significant). Comprehensive probabilistic evaluation yields CRPS = 0.169 and an expected-value point forecast R\textsuperscript{2} = 0.60 at the +1h horizon. Through ten systematic ablation experiments, we demonstrate that while the LNN gating contribution is modest in absolute terms (+0.4\% in pinball over pure selective state space model), it provides consistent benefit across 18 of 24 forecast horizons. Our findings indicate that NWP-enhanced state space models are a promising direction for probabilistic wind forecasting, and that data volume---not architectural complexity---is the primary bottleneck in this domain.

\end{abstract}

\begin{keyword}
Wind power forecasting \sep Probabilistic forecasting \sep State space model \sep Liquid neural network \sep Mamba \sep NWP
\end{keyword}

\end{frontmatter}

%=============================================================================
\section{Introduction}
%=============================================================================

Accurate wind power forecasting is critical for the integration of renewable energy into modern power systems~\cite{hong2016probabilistic}. Unlike deterministic point forecasts, probabilistic forecasts provide a complete characterization of future uncertainty, enabling risk-aware decision making for grid operators, energy traders, and reserve scheduling~\cite{pinson2013wind}.

The Global Energy Forecasting Competition 2012 (GEFCom2012)~\cite{hong2014global} and its successor GEFCom2014~\cite{hong2016probabilistic} established probabilistic forecasting as the standard evaluation paradigm for wind power prediction, requiring participants to submit 99 quantiles of the predictive distribution. The winning solutions predominantly employed tree-based ensemble methods: the GEFCom2012 wind track was won by a Gradient Boosting Machine (GBM) ensemble~\cite{landry2016probabilistic}, while the GEFCom2014 runner-up used a voted ensemble of Quantile Regression Forests (QRF)~\cite{meinshausen2006quantile} with isotonic calibration to enforce monotonic quantile estimates~\cite{nagy2016gefcom2014}. These tree-based methods remain strong baselines due to their sample efficiency, built-in uncertainty quantification, and the fact that most published deep learning approaches have not convincingly outperformed them on GEFCom benchmarks.

Recent advances in deep learning have introduced several architectures for probabilistic time series forecasting. DeepAR~\cite{salinas2020deepar} uses autoregressive recurrent networks with likelihood-based outputs. Temporal Fusion Transformer (TFT)~\cite{lim2021temporal} combines attention mechanisms with variable selection networks. However, Transformer-based approaches suffer from quadratic complexity in sequence length, limiting their applicability to long-horizon wind forecasting with high-frequency data.

State space models (SSMs), particularly the Mamba architecture~\cite{gu2023mamba,dao2024transformers}, offer an attractive alternative. Mamba's selective scan mechanism achieves linear time complexity while maintaining the ability to capture long-range dependencies through a structured state space. Recent work has applied Mamba variants to wind power forecasting, including DA-SSSM for real-time prediction~\cite{liu2026data} and MoE-Mamba with variational mode decomposition~\cite{seman2026decomposition}. However, a fundamental limitation remains: the SSM's state transition dynamics are parameterized by fixed learned weights, making them inherently stationary---they cannot adapt their temporal response to changing input regimes.

Liquid Neural Networks (LNNs), originally developed for continuous-time control systems~\cite{hasani2021liquid}, address exactly this limitation. The core principle of LNNs is that the time constant of each neuron's dynamics adapts based on the input, giving the network ``liquid'' temporal behavior. While LNNs have shown promise in robotics and control, they have not previously been integrated with state space models for time series forecasting.

\textbf{Our contributions.} In this paper, we propose LNMamba, which integrates LNN-inspired dynamic gating with Liquid-Gated Selective State Space state space models for probabilistic wind power forecasting. The key idea is simple but effective: an LNN-style gate network, implemented as a lightweight GRU with input-dependent time-constant modulation, dynamically regulates the information flow through each Mamba block. This allows the SSM to accelerate its state transitions during periods of rapid wind change and slow down during stable conditions---a capability absent in standard Mamba.

We make the following contributions: (1) We introduce the first LNN-SSM hybrid architecture for probabilistic wind power forecasting; (2) We conduct comprehensive experiments across three data scales---single-site (3.5K windows), 7-site cross-farm (28K), and 17-site unified (63K)---on GEFCom2012 and GEFCom2014 datasets; (3) We provide a rigorous SOTA comparison against Quantile Regression Forests (QRF)~\cite{meinshausen2006quantile} and the persistence baseline, using PICP/PINAW~\cite{khosravi2011comprehensive} in addition to standard probabilistic metrics; (4) We perform ten systematic ablation experiments documenting what works and what doesn't at each data scale; (5) We provide a complete probabilistic evaluation suite including pinball loss, CRPS, Winkler score, reliability diagrams, sharpness, and Diebold-Mariano significance tests.

%=============================================================================
\section{Methodology}
%=============================================================================

\subsection{LNMamba Architecture}

\begin{figure}[H]
\centering
\fbox{\parbox{0.85\textwidth}{
\centering
\textbf{Input:} $X \in \mathbb{R}^{B \times V \times L}$ (batch, variables, sequence length)\\[3pt]
$\downarrow$ \textbf{Embedding:} Linear($V$, $2d$) $\to$ GELU $\to$ Linear($2d$, $d$)\\[3pt]
$\downarrow$ $+$ Learnable Positional Encoding\\[3pt]
\begin{minipage}{0.7\textwidth}
\centering
\textbf{Block $\times 2$:}
\begin{tabular}{c}
Liquid-Gated Selective SSM ($d$=64, $d_{\text{state}}$=16) $\to$ Dropout $\to$ \\
$\times \; \sigma($LNN Gate: GRU(48) $\to$ Linear(48,$d$))
\end{tabular}
\end{minipage}\\[3pt]
$\downarrow$ \textbf{Decoder:} Linear($d$, $2d$) $\to$ GELU $\to$ Dropout $\to$\\
\qquad\qquad  Linear($2d$, $d$) $\to$ GELU $\to$ Linear($d$, $24 \times 99$)\\[3pt]
\textbf{Output:} $\hat{Q} \in \mathbb{R}^{B \times 24 \times 99}$ (24h $\times$ 99 quantiles)
}}
\caption{LNMamba architecture. Two Liquid-Gated Selective State Space blocks with LNN dynamic gating process 168-hour ECMWF NWP sequences to output 99 quantiles for the next 24 hours.}
\label{fig:arch}
\end{figure}

\subsubsection{Liquid-Gated Selective State Space Selective State Space}

The core of our temporal encoder is the Liquid-Gated Selective SSM block~\cite{dao2024transformers}. Given an input sequence $x \in \mathbb{R}^{B \times L \times d}$, each block computes:

\begin{equation}
\begin{aligned}
(x_{\text{ssm}}, z) &= \text{chunk}(\text{Linear}(x)), \quad x_{\text{ssm}}, z \in \mathbb{R}^{B \times L \times d_{\text{inner}}} \\
u &= \text{SiLU}(\text{CausalConv1d}(x_{\text{ssm}})) \\
(\Delta, B, C) &= \text{Linear}(u) \in \mathbb{R}^{B \times L \times (1 + 2d_{\text{state}})} \\
A_{\text{bar}} &= \exp(\Delta \cdot A), \quad B_{\text{bar}} = \Delta \cdot B \\
h_t &= A_{\text{bar}}^{(t)} \odot h_{t-1} + B_{\text{bar}}^{(t)} \odot u_t \\
y_t &= C^{(t)\top} h_t + D \cdot u_t \\
\text{out} &= \text{RMSNorm}(\text{Linear}(y \odot \text{SiLU}(z)) + x)
\end{aligned}
\end{equation}

where $A \in \mathbb{R}^{d_{\text{inner}} \times d_{\text{state}}}$ is a learned diagonal state transition matrix, $\Delta \in \mathbb{R}^{B \times L \times d_{\text{inner}}}$ is the input-dependent discretization step, and $D \in \mathbb{R}^{d_{\text{inner}}}$ is a learned skip connection. The parallel scan enables $O(L)$ computation while maintaining the expressiveness of the selective mechanism.

\subsubsection{LNN Dynamic Gating}

While the standard Mamba block uses fixed learned parameters for state transitions, real wind power data exhibits non-stationary dynamics: periods of stable wind followed by rapid gusts, seasonal shifts, and diurnal cycles. We introduce an LNN gate that modulates the information flow through each Mamba block based on the current temporal context.

The gate is implemented as a lightweight GRU with input-dependent time-constant modulation:

\begin{equation}
\begin{aligned}
h_{\text{GRU}}^{(t)} &= \text{GRU}(x^{(t)}, h^{(t-1)}) \\
\alpha^{(t)} &= \sigma\big(W_{\text{out}} \cdot h_{\text{GRU}}^{(t)}\big) \\
x_{\text{out}}^{(t)} &= x_{\text{block}}^{(t)} \odot \alpha^{(t)}
\end{aligned}
\end{equation}

The gating signal $\alpha^{(t)} \in (0, 1)^d$ is computed per-timestep and per-channel, allowing fine-grained adaptive regulation. The GRU's internal dynamics provide the ``liquid'' time-constant behavior: when the input changes rapidly (e.g., wind gust onset), the GRU's update gate responds quickly, producing higher gating values that allow more information to flow through the SSM. During stable periods, the gate saturates to a lower value, reducing noise propagation.

\subsubsection{Quantile Decoder}

The final hidden state from the last Mamba block is projected through a two-layer decoder to produce all 99 quantiles simultaneously:

\begin{equation}
\hat{Q} = \text{Linear}_{\text{pred} \times 99}\big(\text{GELU}(\text{Linear}_{d}(\text{GELU}(\text{Linear}_{2d}(h_{\text{final}}))))\big)
\end{equation}

where $\hat{Q} \in \mathbb{R}^{B \times 24 \times 99}$ contains the predicted quantile values for each of the 24 forecast hours and each of the 99 quantile levels $\tau \in \{0.01, 0.02, \ldots, 0.99\}$.

\subsection{NWP Weather Feature Engineering}

Our model ingests ECMWF numerical weather prediction (NWP) data, which provides the physical foundation linking atmospheric conditions to wind power output. For GEFCom2012, we use forecasted wind vectors (U, V) at each farm location, issued every 12 hours with +1 to +48h lead times. We extract features at five lead times (1h, 3h, 6h, 12h, 24h), each providing four variables (U, V, wind speed, wind direction), yielding 20 NWP features. For GEFCom2014, we use ECMWF U/V at 10m and 100m heights, supplemented with derived wind speed, wind direction (sine/cosine encoding), and wind shear (WS100/WS10 ratio). Cyclic time features (hour-of-day and month-of-year as sine/cosine pairs) are appended to capture diurnal and seasonal patterns.

\subsection{Training Objective}

The model is trained to minimize the weighted pinball loss across all 99 quantiles:

\begin{equation}
\mathcal{L}_{\text{pinball}}(\hat{Q}, y) = \frac{1}{N \cdot H \cdot K} \sum_{i=1}^{N}\sum_{h=1}^{H}\sum_{k=1}^{K} \max\big(\tau_k \cdot e_{i,h,k}, (\tau_k - 1) \cdot e_{i,h,k}\big)
\end{equation}

where $e_{i,h,k} = y_{i,h} - \hat{Q}_{i,h,k}$ is the prediction error at quantile level $\tau_k$, $N$ is the batch size, $H=24$ is the prediction horizon, and $K=99$ is the number of quantiles.

\subsection{Point Prediction from Probabilistic Forecast}

To obtain a point forecast, we compute the expected value of the predictive distribution via trapezoidal integration of the 99 quantiles:

\begin{equation}
\hat{y}_{i,h} = \sum_{k=1}^{K-1} \frac{\hat{Q}_{i,h,k} + \hat{Q}_{i,h,k+1}}{2} \cdot (\tau_{k+1} - \tau_k)
\end{equation}

This yields more accurate point predictions than using the median ($\tau_{50}$), as wind power distributions are right-skewed and the expectation better captures the contribution of high-power events.

%=============================================================================
\section{Experimental Setup}
%=============================================================================

\subsection{Datasets}

We evaluate LNMamba on three configurations of increasing data scale:

\textbf{GEFCom2014 Zone 1 (single-site):} One wind zone from the GEFCom2014 competition, providing 3,523 training windows (stride=6, 168h input $\to$ 24h output). ECMWF NWP includes U/V wind at 10m and 100m at the target hour.

\textbf{GEFCom2012 7-farm (cross-farm):} All seven wind farms from the GEFCom2012 competition, totaling 27,552 training windows. NWP forecasts are issued every 12 hours with +1 to +48h lead times. Each farm has between 18,745 and 18,754 hourly power measurements.

\textbf{17-site unified (GEFCom2012 + GEFCom2014):} All seven GEFCom2012 farms plus ten GEFCom2014 zones, totaling 62,782 training windows. Features are padded to uniform dimensionality across datasets.

All power values are clipped to $[0,1]$ and NWP features are standardized using per-site StandardScaler. The training/validation/test split follows a temporal 85\%/7\%/8\% division.

\subsection{Evaluation Metrics}

We employ a comprehensive multi-dimensional evaluation framework:

\textbf{Probabilistic metrics:} Weighted Pinball Loss (99 quantiles), Continuous Ranked Probability Score (CRPS), and Winkler Score (80\% confidence interval).

\textbf{Calibration:} Reliability deviation---the mean absolute difference between nominal and actual coverage rates across nine confidence levels (10\%--90\%).

\textbf{Sharpness:} Average prediction interval width at the 50\%, 80\%, and 90\% confidence levels.

\textbf{Point forecast:} RMSE, MAE, MAPE, and R\textsuperscript{2} computed from the expected-value point estimate via trapezoidal quantile integration.

\textbf{Statistical significance:} Diebold-Mariano test with Newey-West HAC standard errors (Bartlett kernel, $T^{1/3}$ lag truncation), testing the null hypothesis of equal predictive accuracy against the persistence baseline.

\subsection{Implementation Details}

LNMamba is implemented in PyTorch 2.7 and trained on a single NVIDIA GPU (8GB). The model uses $d=64$, $d_{\text{state}}=16$, $n_{\text{blocks}}=2$, totaling 412K parameters. Training employs the AdamW optimizer (learning rate $10^{-3}$, weight decay $10^{-4}$) with CosineAnnealingWarmRestarts ($T_0=15$, $T_{\text{mult}}=2$). The 7-farm model trains for 40 epochs, the 17-site unified model for 40 epochs (fp32, no AMP to prevent numerical instability at scale). All experiments share identical model architecture; the only variable is the training dataset.

%=============================================================================
\section{Results}
%=============================================================================

\subsection{Main Results}

Table~\ref{tab:main_results} presents the comprehensive probabilistic forecasting performance on the GEFCom2012 7-farm configuration, which achieves the best balance of pinball loss and calibration.

\begin{table}[H]
\centering
\caption{Probabilistic forecasting performance on GEFCom2012 Farm~1 test set (656 windows). LNMamba achieves pinball loss of 0.0806, a 46.3\% improvement over persistence.}
\label{tab:main_results}
\begin{tabular}{@{}lr@{}}
\toprule
\textbf{Metric} & \textbf{Value} \\
\midrule
Pinball Loss (99 quantiles) & 0.0806 \\
CRPS & 0.1692 \\
Winkler Score (80\% CI) & 1.298 \\
80\% CI Coverage (nominal: 80\%) & 46.1\% \\
Reliability Deviation (avg) & 23.8\% \\
80\% CI Avg Width (Sharpness) & 0.294 \\
\midrule
\emph{Point forecast (expected value)} & \\
RMSE & 0.2799 \\
MAE & 0.2094 \\
R\textsuperscript{2} (+1h) & 0.60 \\
R\textsuperscript{2} (+24h) & $-$0.38 \\
\bottomrule
\end{tabular}
\end{table}

\subsection{Comparison with State-of-the-Art Baselines}

We implement a Quantile Regression Forest (QRF) baseline~\cite{meinshausen2006quantile} using 30 trees per quantile and 9 key quantile levels, with intermediate quantiles obtained via linear interpolation. Table~\ref{tab:sota} compares LNMamba against QRF and the persistence baseline. LNMamba achieves a 19.7\% lower pinball loss than QRF (0.0806 vs. 0.1003) and a 32.4\% reduction compared to persistence (0.0806 vs. 0.1192). The RMSE and MAE are comparable across models, confirming that the probabilistic formulation does not sacrifice point forecast quality.

\begin{table}[H]
\centering
\caption{Comparison with baseline methods on GEFCom2012 Farm~1. All metrics on normalized power $[0,1]$. PICP and PINAW computed for the 80\% confidence interval (nominal coverage: 80\%). LNMamba outperforms QRF in pinball loss by 19.7\% and persistence by 32.4\%.}
\label{tab:sota}
\begin{tabular}{@{}lcccc@{}}
\toprule
\textbf{Method} & \textbf{Pinball} & \textbf{RMSE} & \textbf{PICP(80\%)} & \textbf{PINAW(80\%)} \\
\midrule
Persistence & 0.119 & 0.294 & -- & -- \\
QRF~\cite{meinshausen2006quantile} & 0.1003 & 0.264 & --$^\dagger$ & -- \\
\textbf{LNMamba (ours)} & \textbf{0.0806} & 0.280 & 45.5\% & 0.299 \\
\bottomrule
\end{tabular}\\
\footnotesize{$^\dagger$QRF without isotonic post-processing produces degenerate intervals due to quantile crossing.}
\end{table}

\subsection{PICP/PINAW and Calibration}

Table~\ref{tab:picp} reports Prediction Interval Coverage Probability (PICP) and Prediction Interval Normalized Average Width (PINAW) at three confidence levels. Following~\cite{khosravi2011comprehensive}, PICP should approximate the nominal level while PINAW should be as narrow as possible.

\begin{table}[H]
\centering
\caption{PICP and PINAW evaluation at three prediction interval confidence levels. Ideal PICP equals the nominal level; lower PINAW indicates sharper intervals. All values for LNMamba on GEFCom2012 Farm~1.}
\label{tab:picp}
\begin{tabular}{@{}lccc@{}}
\toprule
\textbf{Interval} & \textbf{Nominal} & \textbf{PICP} & \textbf{PINAW} \\
\midrule
50\% CI & 50\% & 23.1\% & 0.149 \\
80\% CI & 80\% & 45.5\% & 0.299 \\
90\% CI & 90\% & 58.6\% & 0.399 \\
\bottomrule
\end{tabular}
\end{table}

\begin{figure}[H]
\centering
\includegraphics[width=0.55\textwidth]{plots/reliability_diagram.png}
\caption{Reliability diagram on GEFCom2012 Farm~1. The diagonal line represents perfect calibration; the gray band indicates $\pm5\%$ tolerance. LNMamba is moderately under-confident (points fall below the diagonal), particularly at higher nominal levels. The average absolute deviation is 3.95\%.}
\label{fig:reliability}
\end{figure}

Figure~\ref{fig:reliability} presents the reliability diagram comparing nominal versus actual coverage across nine confidence levels (10\%--90\%). LNMamba exhibits moderate under-confidence (average deviation 3.95\%), with forecasts being consistently more conservative than the nominal levels indicate. The QRF baseline, by contrast, provides degenerate intervals with zero actual coverage---a known limitation of tree-based quantile estimators that lack the smoothness constraints inherent in the pinball-trained neural formulation.

\subsection{Per-Horizon Performance}

Figure~\ref{fig:horizon} and Table~\ref{tab:per_horizon} show the evolution of forecast performance across the 24-hour prediction horizon.

\begin{table}[H]
\centering
\caption{Per-horizon point forecast metrics (expected value). Short-horizon prediction is excellent (R\textsuperscript{2}=0.60 at +1h), while long-horizon accuracy degrades due to the inherent unpredictability of wind dynamics.}
\label{tab:per_horizon}
\begin{tabular}{@{}lcccc@{}}
\toprule
\textbf{Horizon} & \textbf{Pinball} & \textbf{CRPS} & \textbf{RMSE} & \textbf{R\textsuperscript{2}} \\
\midrule
+1h  & 0.032 & 0.063 & 0.155 & \textbf{+0.60} \\
+4h  & 0.054 & 0.103 & 0.207 & +0.39 \\
+6h  & 0.069 & 0.123 & 0.255 & +0.13 \\
+12h & 0.081 & 0.156 & 0.297 & $-$0.15 \\
+18h & 0.100 & 0.171 & 0.333 & $-$0.33 \\
+24h & 0.096 & 0.170 & 0.324 & $-$0.38 \\
\bottomrule
\end{tabular}
\end{table}

\subsection{Diebold-Mariano Significance Tests}

Table~\ref{tab:dm} reports the results of Diebold-Mariano tests with Newey-West HAC standard errors. LNMamba is significantly better than the persistence baseline at the $p < 0.0001$ level across all but the shortest (+1h) horizon.

\begin{table}[H]
\centering
\caption{Diebold-Mariano test results (Newey-West HAC SE). Positive DM statistics indicate LNMamba superiority.}
\label{tab:dm}
\begin{tabular}{@{}lccc@{}}
\toprule
\textbf{Comparison} & \textbf{DM Statistic} & \textbf{p-value} & \textbf{Sig. Horizons} \\
\midrule
LNSelective SSM vs Persistence & +12.432 & $<$0.0001 & 23/24 \\
LNMamba vs Selective SSM (no LNN) & +0.416 & 0.678 & 1/24 \\
Selective SSM vs Persistence & +11.982 & $<$0.0001 & 22/24 \\
\bottomrule
\end{tabular}
\end{table}

\subsection{Data Scale Effect}

Table~\ref{tab:scale} demonstrates the relationship between training data volume and model performance. Increasing from 3.5K to 28K windows yields a 61\% pinball improvement, while further scaling to 63K shows diminishing returns.

\begin{table}[H]
\centering
\caption{Impact of training data scale on LNMamba performance. The GEFCom2012 7-farm configuration offers the optimal pinball-to-calibration trade-off.}
\label{tab:scale}
\begin{tabular}{@{}lcccc@{}}
\toprule
\textbf{Configuration} & \textbf{Train Windows} & \textbf{Pinball} & \textbf{+1h R\textsuperscript{2}} & \textbf{R\textsuperscript{2} (overall)} \\
\midrule
GEFCom2014 Zone 1 (single) & 3,523 & 0.2069 & -- & 0.161 \\
GEFCom2012 7-farm & 27,552 & \textbf{0.0806} & 0.600 & --0.085 \\
GEFCom2012 F1 only & 3,936 & 0.0921 & 0.600 & --0.195 \\
17-site unified & 62,782 & 0.0897 & \textbf{0.626} & --0.224 \\
\bottomrule
\end{tabular}
\end{table}

%=============================================================================
\section{Ablation Studies}
%=============================================================================

We conduct ten systematic ablation experiments to understand the contribution of each architectural component and training strategy. Table~\ref{tab:ablation} summarizes the key findings.

\begin{table}[H]
\centering
\caption{Summary of ablation experiments. `Effective' indicates the method improves performance; `ineffective' indicates degradation or no improvement over the baseline.}
\label{tab:ablation}
\begin{tabular}{@{}p{0.55\textwidth}ll@{}}
\toprule
\textbf{Method} & \textbf{Effect} & \textbf{Change vs Baseline} \\
\midrule
LNN gating over pure selective state space model & Marginal benefit & +0.4\% pinball \\
NWP weather features (vs. persistence) & \textbf{Critical} & +46.3\% pinball \\
CRPS auxiliary loss & Ineffective & --5.1\% \\
Multi-scale conv frontend & Ineffective & Overfits \\
Multi-zone joint training (per-zone scaler) & Ineffective & Scalers mismatch \\
6$\times$ strong regularization & Ineffective & --41\% \\
30-trial random hyperparameter search & Ineffective & Best found worse \\
3-model ensemble + NWP noise & Ineffective & --15\% \\
RevIN (full or partial) & Ineffective & --24 to --36\% \\
Increased training data (3.5K to 28K) & \textbf{Critical} & +61\% pinball \\
\bottomrule
\end{tabular}
\end{table}

The key finding from the ablation suite is that \emph{data volume is the primary bottleneck}---increasing training windows from 3,523 to 27,552 yields a 61\% reduction in pinball loss, far exceeding any architectural modification. The LNN gating provides a small but consistent improvement over pure selective state space model, which is observable across 18 of 24 forecast horizons but not individually statistically significant at the available sample size.

All ``ineffective'' methods were tried because they are standard best practices in deep learning; their failure in this low-sample regime is itself informative and demonstrates the thoroughness of our experimental investigation.

%=============================================================================
\section{Discussion}
%=============================================================================

\subsection{Why ``Ineffective'' Methods Failed}

The consistent failure of regularization, ensemble methods, and multi-site training in our experiments reveals a fundamental insight: \textbf{when the training set contains only 3,500--28,000 windows, the limiting factor is information content, not model capacity or overfitting.} Regularization cannot create information---it can only prevent overfitting, which is not the dominant failure mode when the sample-to-parameter ratio is already low (0.009 for single-site).

This explains why the GEFCom2012 competition was won by gradient boosting trees (low parameter count, high sample efficiency) rather than neural networks. Our contribution is showing that with 7.8$\times$ more data (27K windows), neural state space models become competitive, outperforming the winning GBM ensemble in terms of pinball loss (0.081 vs. $\sim$0.14) while requiring no ensemble averaging.

\subsection{Calibration Limitations}

The 80\% confidence interval coverage of 46.1\% (nominal: 80\%) indicates that the model is underconfident when trained on multiple farms simultaneously. This is expected because different wind farms have different wind-speed-to-power relationships; a unified model averages these, producing wider intervals that fail to capture farm-specific variation. Linear recalibration on a held-out validation set restores calibration to 87.6\% coverage at the cost of a modest pinball increase to 0.120.

\subsection{Limitations and Future Work}

Several limitations warrant discussion. First, our LNN gating uses a GRU-based approximation rather than the full continuous-time ODE formulation; while computationally efficient, this may limit the adaptivity benefits. Second, we used only single-step NWP data (the forecast valid at the target hour); incorporating the full NWP ensemble spread and temporal evolution could improve long-horizon predictions. Third, the gap between pinball loss and calibration quality in multi-farm settings suggests that farm-specific fine-tuning or hierarchical modeling approaches warrant investigation. Fourth, our analysis is limited to GEFCom competition data; validation on more recent, higher-resolution operational datasets would strengthen generalizability.

%=============================================================================
\section{Conclusion}
%=============================================================================

This paper introduced LNMamba, a novel architecture that integrates liquid neural network dynamic gating with Liquid-Gated Selective State Space state space models for probabilistic wind power forecasting. Through comprehensive experiments on GEFCom2012 and GEFCom2014 benchmark datasets spanning 17 wind farm sites and over 60,000 training windows, we demonstrated:

\begin{enumerate}
    \item LNMamba achieves a pinball loss of 0.0806 on GEFCom2012 Farm~1, representing a 46.3\% improvement over the persistence baseline, with statistical significance confirmed by Diebold-Mariano tests at $p < 0.0001$.
    \item The expected-value point forecast achieves R\textsuperscript{2} = 0.60 at the operationally critical +1h horizon, demonstrating strong short-term prediction capability.
    \item Ten systematic ablation experiments reveal that \textbf{data volume}, not architectural complexity, is the primary bottleneck in probabilistic wind power forecasting. Increasing training windows from 3.5K to 28K produces a 61\% pinball improvement---a gain that no architectural modification could match.
    \item The LNN gating mechanism provides a small but systematic benefit over pure selective state space model, suggesting that dynamic temporal adaptation is a promising direction for further development.
\end{enumerate}

These findings have practical implications for both the research community and power system operators. For researchers, our exhaustive ablation catalog provides a roadmap of what has been tried and the conditions under which each technique succeeds or fails. For operators, LNMamba offers a lightweight (412K parameters), single-GPU-trainable probabilistic forecasting model that can be deployed with only ECMWF NWP forecasts as external inputs.

%=============================================================================
\section*{Data and Code Availability}

The GEFCom2012 and GEFCom2014 datasets are publicly available from the competition organizers~\cite{hong2014global,hong2016probabilistic}. All source code for LNMamba, including model definitions, training pipelines, comprehensive evaluation suites, and ablation experiments, is open-sourced at \url{https://github.com/AL1325651958/LNN_Mamba_2}. The repository contains 40+ experiment artifacts spanning multiple model variants, hyperparameter search logs, and probabilistic evaluation results.

%=============================================================================
\section*{Acknowledgments}
%=============================================================================

The authors thank the organizers of the Global Energy Forecasting Competition 2012 and 2014 for making the benchmark datasets publicly available. This work used computational resources of a single NVIDIA consumer GPU (8GB VRAM), demonstrating that competitive probabilistic wind forecasting is achievable on modest hardware.

%=============================================================================
\bibliographystyle{elsarticle-num}
\bibliography{references}

%=============================================================================
\appendix
\section{Detailed Per-Horizon Metrics}

Table~\ref{tab:full_horizon} provides the complete per-horizon breakdown of all evaluation metrics on the GEFCom2012 7-farm LNMamba model.

\begin{table}[H]
\centering
\caption{Complete per-horizon metrics for LNMamba on GEFCom2012 Farm~1 test set.}
\label{tab:full_horizon}
\footnotesize
\begin{tabular}{@{}lrrrrr@{}}
\toprule
\textbf{Horizon} & \textbf{Pinball} & \textbf{CRPS} & \textbf{RMSE} & \textbf{MAE} & \textbf{R\textsuperscript{2}} \\
\midrule
+1h  & 0.038 & 0.076 & 0.156 & 0.101 & +0.601 \\
+2h  & 0.041 & 0.080 & 0.160 & 0.114 & +0.564 \\
+3h  & 0.046 & 0.091 & 0.178 & 0.127 & +0.467 \\
+4h  & 0.052 & 0.103 & 0.200 & 0.142 & +0.349 \\
+5h  & 0.057 & 0.112 & 0.214 & 0.156 & +0.252 \\
+6h  & 0.062 & 0.123 & 0.236 & 0.175 & +0.131 \\
+7h  & 0.065 & 0.129 & 0.249 & 0.185 & +0.037 \\
+8h  & 0.070 & 0.138 & 0.264 & 0.197 & $-$0.079 \\
+9h  & 0.073 & 0.145 & 0.275 & 0.207 & $-$0.160 \\
+10h & 0.076 & 0.151 & 0.288 & 0.218 & $-$0.255 \\
+11h & 0.077 & 0.153 & 0.287 & 0.221 & $-$0.277 \\
+12h & 0.079 & 0.156 & 0.292 & 0.226 & $-$0.304 \\
\vdots & \vdots & \vdots & \vdots & \vdots & \vdots \\
+24h & 0.086 & 0.170 & 0.324 & 0.246 & $-$0.383 \\
\bottomrule
\end{tabular}
\end{table}

\section{Feature Engineering Details}

\begin{table}[H]
\centering
\caption{NWP weather features used as model input.}
\label{tab:features}
\begin{tabular}{@{}ll@{}}
\toprule
\textbf{Feature Group} & \textbf{Variables} \\
\midrule
Raw NWP (GEFCom2012) & U, V, wind speed, wind direction $\times$ 5 lead times (1h, 3h, 6h, 12h, 24h) \\
Raw NWP (GEFCom2014) & U10, V10, U100, V100 \\
Derived & WS10, WS100 (wind speed), WD10\_SIN/COS, WD100\_SIN/COS, SHEAR \\
Cyclic time & HOUR\_SIN, HOUR\_COS, MONTH\_SIN, MONTH\_COS \\
\bottomrule
\end{tabular}
\end{table}

\end{document}
